% ===================================================================
%  CADRO N.º 2 — O corpo humano. — O recreo. — Os xogos.
%  Traduction 2026 d'apres texto/io, controlee sur texto/fr, texto/es
%  et texto/pt. Consigne : voir texto/gl/10-cadro-01.tex.
%
%  LE TABLEAU DU CORPS EST LE MEILLEUR BANC D'ESSAI DE CETTE COLONNE,
%  parce que le corps s'apprend a la maison et dans aucune ecole. Le
%  galicien y separe donc ses deux voisines mot pour mot, et sans
%  aucun emprunt : « meixela » la joue, ou l'espagnol dit « mejilla »
%  et le portugais « bochecha » ; « beizo » la levre, contre « labio »
%  et « lábio » ; « xeonllo » le genou, contre « rodilla » et
%  « joelho » ; « cóbado » le coude, contre « codo » et « cotovelo » ;
%  « nocello » la cheville, contre « tobillo » et « tornozelo » ;
%  « pescozo » le cou, contre « cuello » et « pescoço ». Sur ces six
%  mots, le galicien ne suit ni l'une ni l'autre quatre fois.
%
%  ET C'EST LA QUE LA COPIE SE VERRAIT. Une colonne galicienne ecrite
%  sur l'espagnol dirait « rodilla » ; ecrite sur le portugais, elle
%  dirait « joelho » avec un J que le galicien n'ecrit pas. Le
%  balayage attrape le second cas et pas le premier — d'ou la
%  relecture, qui reste le seul controle du premier.
%
%  LES JEUX, EUX, N'ONT PAS DE NOM INTERNATIONAL, et l'ido le dit
%  lui-meme dans la note du feuillet 17 : il a du les nommer par
%  l'action, faute d'accord entre les langues. Le galicien a les
%  siens, qui ne sont ni ceux de Madrid ni ceux de Lisbonne : « a pita
%  cega » pour le colin-maillard, « as agochadas » pour cache-cache,
%  « a man quente », « o xogo da ra », « a randeeira ». On les ecrit
%  tels quels ; les traduire depuis l'espagnol aurait donne « la
%  gallinita ciega » deguisee, ce qui n'est le jeu de personne.
% ===================================================================

%%K t02-apar-1 apar
\VUsaut{4.0mm}
\VUtitre{15.0pt}{60}{CADRO N.\textsuperscript{o} 2}
\VUsaut{2.0mm}
\VUtitre{11.4pt}{0}{O corpo humano. --- O recreo.}
\VUtitre{11.4pt}{0}{Os xogos.}

%%K t02-tit-0 sub
\VUtitre{13.2pt}{0}{O corpo humano.}
\VUsaut{2.4mm}
\VUcentre{10.2pt}{0}{\textit{(Lección sinxela de ciencias naturais.)}}
\VUsaut{3.0mm}

%%K t02-01-1 p
1. --- As principais partes do corpo humano son: a \VUgras{cabeza,} o \VUgras{tronco}
e os \VUgras{membros.}

%%K t02-tit-1 sub
\VUpk{10.2pt}{40}{A cabeza.}
\VUsaut{3.0mm}

%%K t02-02-1 p
2. --- A cabeza comprende dúas partes: o \VUgras{cranio,} que contén o
\VUgras{cerebro} e que a \VUgras{cabeleira} cobre, e a \VUgras{cara.}

%%K t02-03-1 p
3. --- No noso debuxo non vemos nin o cranio nin o cerebro; pero vemos moi claramente
as partes importantes da cara.

%%K t02-04-1 p
4. --- Velaquí primeiro a \VUgras{fronte,} que denota unha intelixencia poderosa,
un pensamento sempre activo. As \VUgras{tempas} están moi despexadas. O \VUgras{ollo}
é vivo e ben aberto. Unha \VUgras{cella} mesta e unhas \VUgras{pestanas} longas
resgárdano.

%%K t02-05-1 p
5. --- Velaquí ademais o \VUgras{nariz} e as \VUgras{ventas.} O nariz sérvenos para
cheirar e para respirar. A cada lado do nariz están as \VUgras{meixelas,} e máis
lonxe as \VUgras{orellas.} A parte baixa da cara consta da \VUgras{boca} e do
\VUgras{queixo.}

%%K t02-06-1 p
6. --- Non os vemos moi ben, pois o \VUgras{bigote} e as \VUgras{patillas}
agóchannolos. Pero sabemos que a boca ten os dous \VUgras{beizos,} o \VUgras{maxilar}
e a \VUgras{mandíbula} cos seus \VUgras{dentes,} e a \VUgras{lingua.} Coa boca
falamos e comemos. Coa lingua saboreamos os alimentos.

%%K t02-07-1 p
7. --- O ollo, a orella, o nariz e a boca son, igual ca os dedos, os órganos dos
cinco sentidos: a \VUgras{vista,} o \VUgras{oído,} o \VUgras{olfacto,} o
\VUgras{gusto,} o \VUgras{tacto.}

%%K t02-08-1 p
8. --- A cabeza é, xa que logo, unha das partes máis interesantes do corpo humano.

%%K t02-tit-2 sub
\VUsaut{4.4mm}
\VUpk{10.2pt}{40}{O tronco.}
\VUsaut{3.0mm}

%%K t02-09-1 p
9. --- O \VUgras{pescozo} xunta a cabeza ao tronco. Comprende a \VUgras{gorxa} e a
\VUgras{caluga.}

%%K t02-10-1 p
10. --- O tronco comprende tamén moitos órganos esenciais. No \VUgras{peito} están
os \VUgras{pulmóns,} cos que respiramos, e o \VUgras{corazón,} que é o centro da
vida. Cando o corazón deixa de bater, o ser vivo morre.

%%K t02-11-1 p
11. --- As \VUgras{costelas} sosteñen o peito.

%%K t02-12-1 p
12. --- As demais partes do tronco son o \VUgras{costado,} as \VUgras{costas,} os
\VUgras{ombreiros} e o \VUgras{abdome} (ventre). Deitámonos sobre o costado ou sobre
as costas para durmir, e levamos as cargas sobre os nosos ombreiros.

%%K t02-13-1 p
13. --- No abdome están o \VUgras{estómago} e os \VUgras{intestinos.} Sérvennos para
dixerir os alimentos.

%%K t02-tit-3 sub
\VUsaut{4.4mm}
\VUpk{10.2pt}{40}{Os membros.}
\VUsaut{3.0mm}

%%K t02-14-1 p
14. --- Temos catro \VUgras{membros}: dous \VUgras{brazos} e dúas \VUgras{pernas.}
Cos brazos collemos, suxeitamos, erguemos e recollemos os obxectos; coas pernas estamos de
pé, camiñamos e corremos.

%%K t02-15-1 p
15. --- O \VUgras{ombreiro} xunta o brazo ao corpo. Un \VUgras{músculo} moi robusto,
o bíceps, move o brazo, que o \VUgras{cóbado} xunta ao \VUgras{antebrazo.} O
\VUgras{carpo} forma a articulación da \VUgras{man,} que remata nos \VUgras{dedos} e
nas \VUgras{uñas.} Na man distinguimos unha \VUgras{vea} (e unha arteria) na que corre
o \VUgras{sangue.}

%%K t02-16-1 p
16. --- As partes da perna son: a \VUgras{coxa,} o \VUgras{xeonllo} e o
\VUgras{xarrete,} a \VUgras{pantorrilla,} cuxa \VUgras{carne} está cuberta pola
\VUgras{pel,} o \VUgras{nocello} e o \VUgras{pé.} Os \VUgras{ósos} da perna son moi
grosos e moi fortes. Na punta do pé están os \VUgras{dedos do pé.} As demais partes
son a \VUgras{planta} e o \VUgras{talón.}

%%K t02-17-1 p
17. --- Tal é a descrición case completa do corpo humano.

%%K t02-17-2 p
\textit{Nota.} --- \textit{Coa axuda da expresión: « dóeme o... (a cabeza, etc.) », o
mestre poderá reempregar todo este vocabulario desde o punto de vista do bo ou mal}
\VUgras{estado corporal.}

%%K t02-tit-4 sub
\VUsaut{5.0mm}
\VUpk{10.2pt}{40}{Ximnasia.}
\VUsaut{2.4mm}
\VUcentre{10.2pt}{0}{\textit{(Relato dun dos alumnos.)}}
\VUsaut{3.0mm}

%%K t02-18-1 p
18. --- Para sermos flexibles, áxiles e sans, temos que exercitar as pernas e os
brazos. Por iso xogamos e practicamos \VUgras{ximnasia.}

%%K t02-19-1 p
19. --- Durante a semana, estes dous exercicios fanse no patio do instituto (véxase o
cadro n.º 1). Pero os xoves e os domingos imos a un gran prado, onde temos sitio para
correr e divertirnos.

%%K t02-20-1 p
20. --- Os nosos mestres e os nosos pais acompáñannos ás veces alí; pero é ben o
\VUgras{mestre de ximnasia} \textsuperscript{(12)} quen dirixe os exercicios.

%%K t02-21-1 p
21. --- No prado, á esquerda da entrada, están os \VUgras{aparellos de ximnasia}
\textsuperscript{(1)}; subimos pola \VUgras{escada} \textsuperscript{(2)}, poñémonos
do revés nas \VUgras{anelas} \textsuperscript{(3)} e no \VUgras{trapecio}
\textsuperscript{(4)}, ízamonos coas mans ata o alto da \VUgras{escada de corda}
\textsuperscript{(5)} ou da \VUgras{corda con nós} \textsuperscript{(6)}.

%%K t02-22-1 p
22. --- Mentres tanto, os nosos compañeiros saltan por riba do \VUgras{cabalo de
madeira} \textsuperscript{(7)} ou das \VUgras{barras paralelas}
\textsuperscript{(11)}. Outros \VUgras{boxean} \textsuperscript{(8)} ou
\VUgras{esgriman} \textsuperscript{(9)} con careta e con \VUgras{florete.} Finalmente,
outros \VUgras{levantan halteres} \textsuperscript{(10)}.

%%K t02-tit-5 sub
\VUsaut{4.6mm}
\VUpk{10.2pt}{40}{Os xogos.}
\VUsaut{3.0mm}

%%K t02-23-1 p
23. --- Arredor dos ximnastas, rapaces e rapazas xogan a \VUgras{xogos} diversos. As
rapazas divírtense co seu \VUgras{aro} \textsuperscript{(14)}; saltan á
\VUgras{corda} \textsuperscript{(13)} ou convidan os irmáns e mais os curmáns a unha
partida de \VUgras{croquet} \textsuperscript{(15)}. Encántalles empurrar lonxe a
\VUgras{bóla} do contrario co seu \VUgras{mazo de madeira} no intre en que el pensa
pasar doadamente o arquiño!

%%K t02-24-1 p
24. --- Os rapaces prefiren os xogos máis musculares. De boa gana xogan aos
\VUgras{birlos} \textsuperscript{(16)}, que derruban con destreza cunha \VUgras{bóla}
\textsuperscript{(17)}. Tampouco desprezan xogar co \VUgras{peón}
\textsuperscript{(18)} e co \VUgras{boliche} \textsuperscript{(19)} ou mesmo ao
\VUgras{xogo da ra} \textsuperscript{(20)}. Pero os seus xogos preferidos son ben as
\VUgras{canicas} \textsuperscript{(21)} e a \VUgras{pelota contra a parede}
\textsuperscript{(22)}, pois o muro do \VUgras{invernadoiro} \textsuperscript{(26)} é
moi axeitado para facer rebotar a pelota lixeira.

%%K t02-25-1 p
25. --- O pequeno Xoán, que anda cos seus \VUgras{zancos} \textsuperscript{(24)}, é
moi hábil e sabe moi ben gardar o equilibrio. Preto del, algúns compañeiros xogan ás
\VUgras{agochadas} \textsuperscript{(25)} nese invernadoiro e detrás da \VUgras{sebe.}
Outros prefiren \VUgras{a man quente} \textsuperscript{(28)} ou o \VUgras{volante}
\textsuperscript{(29)}, que voa dunha \VUgras{raqueta} á outra.

%%K t02-noto-1 noto
\VUnotes{143.2mm}{%
(*) Os xogos de nenos ou de rapaces teñen a miúdo nomes puramente nacionais.
Tratamos de nomealos en Ido o mellor posible, xa inspirándonos na propia acción que
os caracteriza, xa escollendo o nome nacional deste ou daquel país, cando non é
inxusto de máis. Ex.: o \VUgras{xogo do barril} (en alemán e en francés) é o
\VUgras{xogo da ra} \textsuperscript{(20)} (en inglés), no cal os xogadores tentan
guindar os seus discos redondos pola boca da ra de metal que está de pé, coa boca
aberta, sobre a caixa (o barril) furada. O \VUgras{salto ao ombreiro}
\textsuperscript{(30)} difire do \VUgras{salto ás costas} só nisto: para saltaren por
riba da cabeza, os xogadores apoian as mans nos ombreiros do compañeiro, mentres que
no «\VUgras{salto ás costas}» apoian as mans só nas costas del. --- Ou tamén algúns
dos participantes están inclinados mentres os outros saltarán por riba das costas
deles apoiando as mans no ombreiro; os primeiros deben aturar o choque sen dobrarse,
os segundos deben saltar sen perderen o equilibrio (Véxase o propio cadro).%
}

%%K t02-26-1 p
26. --- No medio do prado hai dúas árbores. Ao pé dunha delas, sete ou oito rapaces
organizaron unha partida de \VUgras{salto ao ombreiro} \textsuperscript{(30)} (*).
Non en todos os países se coñece ese xogo; pero todos saben o que é o \VUgras{salto ás
costas,} ao que os rapaces non xogan desta vez.

%%K t02-tit-6 sub
\VUsaut{5.0mm}
\VUpk{10.2pt}{40}{Os xogos ao aire libre.}
\VUsaut{3.4mm}

%%K t02-27-1 p
27. \VUgras{A pita cega} \textsuperscript{(31)} é tamén moi divertida; rapazas e
rapaces poñen á quenda a \VUgras{venda} nos ollos e apañan os torpes que se deixan
coller.

%%K t02-28-1 p
28. --- No medio do prado, velaquí un xogo moi francés pero algo violento: é
\VUgras{o oso} \textsuperscript{(32)}, moito máis ruidoso ca o tan tranquilo xogo do
\VUgras{papaventos} \textsuperscript{(33)} ou mesmo ca o xogo inglés do
\VUgras{tenis} \textsuperscript{(34)}.

%%K t02-29-1 p
29. --- Este último deporte é a ledicia dos alumnos maiores. Os propios mestres practícano de boa
gana. Esixe grande destreza xunto con axilidade e práceme moito.
Déixolles ás rapazas os xogos da \VUgras{randeeira} \textsuperscript{(35)}.

%%K t02-30-1 p
30. --- Non me gusta outro xogo inglés que quizais chegue a ser perigoso.

%%K t02-30-1b p
Quero falar do \VUgras{fútbol} \textsuperscript{(36)}, que alegra o meu irmán maior.
Prefiro o noso vello \VUgras{xogo das barras} \textsuperscript{(38)}, tan hixiénico e
de verdade menos brutal.

%%K t02-tit-7 sub
\VUsaut{4.6mm}
\VUpk{10.2pt}{40}{Ciclismo e globo.}
\VUsaut{3.0mm}

%%K t02-31-1 p
31. --- Tamén de boa gana dou a volta ao prado en \VUgras{bicicleta}
\textsuperscript{(37)}.

%%K t02-32-1 p
32. --- O ano que vén aprenderei a \VUgras{montar a cabalo} \textsuperscript{(39)}.
Que fermoso é un cabalo que anda ou trota con garbo, e que salta a galope por riba
dos obstáculos!

%%K t02-33-1 p
33. --- No verán aprendemos a \VUgras{nadar} \textsuperscript{(40)}. Somerxémonos e
bañámonos con deleite na auga clara e fresca do río. Ás veces, para rematar,
lanzamos un \VUgras{globo} de papel \textsuperscript{(41)} que ascende maxestoso á
atmosfera axiña que está cheo de aire quente.

%%K t02-34-1 p
34. --- Hai tamén moitos outros xogos, pero non acabaría de os enumerar, e
compréndese, segundo este resumo, que os xoves non nos aburrimos moito no noso fermoso
prado.

%%K t02-34-2 p
N.-B. --- \textit{O mestre completará doadamente este capítulo, describindo máis polo
miúdo cada un dos xogos máis importantes.}
